% =====================================================================
% Memorial de cálculo - Trabalho de VE - Mecânica II - Grupo I (Pião)
% Desenvolve, passo a passo, as contas do relatório.
% Compilação: compilar ANTES relatorio_grupo_I.tex (os números das equações
% do relatório são lidos de relatorio_grupo_I.aux) e depois executar
% pdflatex memorial_de_calculo.tex  duas vezes.
% =====================================================================
\documentclass[11pt,a4paper]{article}

\usepackage[utf8]{inputenc}
\usepackage[T1]{fontenc}
\usepackage{lmodern}
\usepackage[brazil]{babel}
\usepackage[a4paper,top=2.5cm,bottom=2.5cm,left=2.5cm,right=2.5cm]{geometry}
\usepackage{amsmath,amssymb}
\usepackage{icomma}
\usepackage{graphicx}
\usepackage{float}
\usepackage{booktabs}
\usepackage{setspace}
\usepackage{indentfirst}
\usepackage{tikz}
\usetikzlibrary{arrows.meta,calc}
\definecolor{azul}{HTML}{2A78D6}
\definecolor{laranja}{HTML}{EB6834}
\definecolor{tinta2}{HTML}{52514E}
\definecolor{apoio}{HTML}{E3E1DA}
\definecolor{corpo}{HTML}{F4E6DC}
\definecolor{corpoesc}{HTML}{DDBBA6}
\definecolor{corpob}{HTML}{A9765A}
\usepackage[font=small,labelfont=bf,labelsep=endash]{caption}
\usepackage{xr-hyper}
\usepackage[hidelinks]{hyperref}
\externaldocument[rel-]{relatorio_grupo_I}

\setstretch{1.15}

\newcommand{\vet}[1]{\mathbf{#1}}
\newcommand{\vomega}{\boldsymbol{\omega}}
\newcommand{\vOmega}{\boldsymbol{\Omega}}
\newcommand{\seno}{\operatorname{sen}}
\newcommand{\un}[1]{\,\mathrm{#1}}
\newcommand{\phid}{\dot{\phi}}
\newcommand{\psid}{\dot{\psi}}
\newcommand{\eps}{\varepsilon}
% remissão à equação do relatório
\newcommand{\rel}[1]{\textit{(equação~\ref*{rel-eq:#1} do relatório)}}

\title{\textbf{Memorial de cálculo}\\[0.2cm]\large Pião: precessão permanente -- Trabalho de VE de Mecânica II -- Grupo I}
\author{}
\date{Novembro de 2026}

\begin{document}
\maketitle
\vspace{-1cm}

\noindent Este documento abre, passo a passo, as contas que o relatório apresenta de forma resumida. Nenhum resultado novo é introduzido: cada dedução parte dos dados do enunciado e termina em uma equação ou valor do relatório, indicado entre parênteses e em itálico. Os valores numéricos intermediários aparecem arredondados, mas os resultados foram calculados sem arredondamentos intermediários.

\tableofcontents

% ---------------------------------------------------------------------
\section{Dados e conversões de unidades}

Todas as grandezas são convertidas para o Sistema Internacional antes de qualquer conta:
\begin{align}
m &= 100\un{g} = 0,100\un{kg}, & c &= 40\un{mm} = 0,040\un{m},\\
k &= 25\un{mm} = 0,025\un{m}, & k' &= 50\un{mm} = 0,050\un{m}.
\end{align}

O \textit{spin} é dado em rotações por minuto. Como uma rotação equivale a $2\pi$ radianos e um minuto a $60$ segundos,
\begin{equation}
\psid = 1800\,\frac{\mathrm{rot}}{\mathrm{min}}\cdot\frac{2\pi\un{rad}}{1\un{rot}}\cdot\frac{1\un{min}}{60\un{s}} = 60\pi\un{rad/s} = 188,4956\un{rad/s}.
\end{equation}
A conversão inversa, usada para expressar os resultados em rpm, é $1\un{rad/s}=60/(2\pi)\un{rpm}=9,5493\un{rpm}$. O enunciado não fornece a gravidade; adota-se $g=9,81\un{m/s^2}$.

% ---------------------------------------------------------------------
\section{Grandezas auxiliares}

O momento de inércia em relação a um eixo é a massa multiplicada pelo quadrado do raio de giração correspondente \rel{inercia}:
\begin{align}
I  &= mk^2  = 0,100\cdot(0,025)^2 = 0,100\cdot6,25\times10^{-4} = 6,25\times10^{-5}\un{kg\cdot m^2},\\
I' &= mk'^2 = 0,100\cdot(0,050)^2 = 0,100\cdot2,50\times10^{-3} = 2,50\times10^{-4}\un{kg\cdot m^2}.
\end{align}
A diferença e os demais produtos que aparecem repetidamente \rel{coef} valem:
\begin{align}
I'-I &= 2,50\times10^{-4}-0,625\times10^{-4} = 1,875\times10^{-4}\un{kg\cdot m^2},\\
W &= mg = 0,100\cdot9,81 = 0,981\un{N},\\
Wc &= 0,981\cdot0,040 = 0,03924\un{N\cdot m},\\
b &= I\psid = 6,25\times10^{-5}\cdot188,4956 = 0,0117810\un{kg\cdot m^2/s}.\label{m:b}
\end{align}

% ---------------------------------------------------------------------
\section{Cinemática}

\begin{figure}[H]
\centering
% Figura: referencial Oxyz (a) e diagrama de corpo livre (b).
% theta = 32 graus (esquemático): \S = sen(theta), \C = cos(theta).
\def\S{0.5299}
\def\C{0.8480}
\def\perfilpiao{%
  plot[smooth,tension=0.6] coordinates {(-0.72,3.1) (-1.0,2.85) (-1.05,2.3) (-0.85,1.6)
    (-0.5,0.85) (-0.18,0.25) (0,0) (0.18,0.25) (0.5,0.85) (0.85,1.6) (1.05,2.3) (1.0,2.85) (0.72,3.1)}}%
\begin{tikzpicture}[scale=0.95,baseline=(cap.base),
    x={(0.966cm,0.129cm)},y={(-0.259cm,0.483cm)},z={(0cm,0.866cm)},
    eixo/.style={-{Stealth[length=2.2mm]},tinta2,line width=0.5pt},
    taxa/.style={-{Stealth[length=2mm]},laranja,line width=1pt}]
  % eixo x: abaixo do plano de apoio, portanto desenhado antes dele
  \draw[eixo] (0,0,0) -- ({2.9*\C},0,{-2.9*\S}) node[right,text=black] {$x\ (\mathbf{i})$};
  % plano de apoio translúcido
  \fill[apoio,opacity=0.8] (-2,-2,0) -- (2,-2,0) -- (2,2,0) -- (-2,2,0) -- cycle;
  \draw[tinta2,line width=0.3pt] (-2,-2,0) -- (2,-2,0) -- (2,2,0) -- (-2,2,0) -- cycle;
  % corpo: silhueta alinhada à projeção do eixo z
  \begin{scope}[x=1cm,y=1cm,rotate=-32.6,yscale=0.95]
    \shade[left color=corpoesc,right color=corpoesc,middle color=corpo,draw=corpob,line width=0.5pt]
      \perfilpiao -- cycle;
    \draw[fill=corpo,draw=corpob,line width=0.5pt] (0,3.1) ellipse (0.72 and 0.22);
  \end{scope}
  % eixos
  \draw[eixo] (0,0,0) -- (0,0,4.8) node[above left,text=black,inner sep=1pt] {$Z\ (\mathbf{K})$};
  \draw[eixo] (0,0,0) -- ({4.9*\S},0,{4.9*\C}) node[right,text=black] {$z\ (\mathbf{k})$};
  \draw[eixo] (0,0,0) -- (0,2.2,0) node[left,text=black,inner sep=1.5pt] {$y\ (\mathbf{j})$};
  % theta, no plano Zz
  \draw[line width=0.5pt] plot[domain=0:32,samples=20] ({1.5*sin(\x)},0,{1.5*cos(\x)});
  \node[inner sep=1pt] at ({1.85*sin(16)},0,{1.85*cos(16)}) {$\theta$};
  % trajetória de G, com o sentido do movimento
  \draw[azul,densely dashed,line width=0.6pt] plot[domain=0:360,samples=90]
    ({2.4*\S*cos(\x)},{2.4*\S*sin(\x)},{2.4*\C});
  \draw[-{Stealth[length=2mm]},azul,line width=0.8pt] plot[domain=20:50,samples=12]
    ({2.4*\S*cos(\x)},{2.4*\S*sin(\x)},{2.4*\C});
  \fill ({2.4*\S},0,{2.4*\C}) circle (1.5pt) node[below right,inner sep=2pt] {$G$};
  \fill (0,0,0) circle (1.5pt) node[below left,inner sep=2pt] {$O$};
  % taxas no sentido positivo (anti-horário visto da ponta de cada eixo)
  \draw[taxa] plot[domain=200:490,samples=50] ({0.42*cos(\x)},{0.42*sin(\x)},4.15);
  \node[left=6pt] at (0,0,4.15) {$\dot{\phi}$};
  \draw[taxa] plot[domain=200:490,samples=50]
    ({4.25*\S+0.42*cos(\x)*\C},{0.42*sin(\x)},{4.25*\C-0.42*cos(\x)*\S});
  \node[right=7pt] at ({4.25*\S},0,{4.25*\C}) {$\dot{\psi}$};
  \node[anchor=north] (cap) at (canvas cs:x=0.35cm,y=-1.45cm) {\small (a) Referencial e movimento};
\end{tikzpicture}\hfill
\begin{tikzpicture}[scale=0.95,baseline=(cap.base),
    eixo/.style={-{Stealth[length=2.2mm]},tinta2,line width=0.5pt},
    forca/.style={-{Stealth[length=2.6mm]},line width=1.3pt}]
  % corpo em vista lateral
  \begin{scope}[rotate=-32]
    \shade[left color=corpoesc,right color=corpoesc,middle color=corpo,draw=corpob,line width=0.5pt]
      \perfilpiao -- cycle;
    \draw[corpob,line width=0.5pt] (-0.72,3.1) -- (0.72,3.1);
  \end{scope}
  % eixos
  \draw[eixo] (0,0) -- (0,4.7) node[above,text=black] {$Z$};
  \draw[eixo] (0,0) -- ({4.9*\S},{4.9*\C}) node[right,text=black] {$z$};
  \draw[eixo] (0,0) -- ({2.1*\C},{-2.1*\S}) node[right,text=black] {$x$};
  % theta
  \draw[line width=0.5pt] (0,1.65) arc[start angle=90,end angle=58,radius=1.65];
  \node[inner sep=1pt] at (0.51,1.92) {$\theta$};
  % cota c, paralela ao eixo z
  \coordinate (A) at ({-1.55*\C},{1.55*\S});
  \coordinate (B) at ({-1.55*\C+2.4*\S},{1.55*\S+2.4*\C});
  \draw[tinta2,densely dotted,line width=0.4pt] (0,0) -- (A);
  \draw[tinta2,densely dotted,line width=0.4pt] ({2.4*\S},{2.4*\C}) -- (B);
  \draw[{Stealth[length=1.8mm]}-{Stealth[length=1.8mm]},line width=0.4pt] (A) -- (B)
    node[midway,above left,inner sep=1pt] {$c$};
  % forças
  \draw[forca,laranja] ({2.4*\S},{2.4*\C}) -- ++(0,-1.6) node[right,text=black,inner sep=2pt] {$\mathbf{W}$};
  \draw[forca,azul] (0,0) -- (0,1.05) node[left,text=black,inner sep=2pt] {$\mathbf{N}$};
  \draw[forca,azul] (0,0) -- (-1.05,0) node[below,text=black,inner sep=2pt] {$\mathbf{R}_h$};
  \fill ({2.4*\S},{2.4*\C}) circle (1.5pt) node[above right,inner sep=2pt] {$G$};
  % O com o eixo y entrando no plano (símbolo: círculo com X)
  \draw[fill=white,line width=0.5pt] (0,0) circle (2.6pt);
  \draw[line width=0.5pt] (-1.84pt,-1.84pt) -- (1.84pt,1.84pt) (-1.84pt,1.84pt) -- (1.84pt,-1.84pt);
  \node[below,inner sep=4pt] at (0,0) {$O$};
  \node[anchor=north,align=center,font=\footnotesize] at (-0.35,-0.62) {$y$ entra no plano};
  \node[anchor=north] (cap) at (0.45,-1.45) {\small (b) Diagrama de corpo livre};
\end{tikzpicture}

\caption{Referencial $Oxyz$ e forças no pião (a mesma figura do relatório).}
\label{m:fig:dcl}
\end{figure}

A Figura~\ref{m:fig:dcl} define o referencial. O eixo $z$ é o eixo de simetria; $x$ pertence ao plano vertical $Zz$; $y$ completa o triedro direto, com $\vet{j}=\vet{k}\times\vet{i}$.

\paragraph{Decomposição de $\vet{K}$.} O vetor unitário vertical pertence ao plano $xz$, logo só tem componentes segundo $\vet{i}$ e $\vet{k}$. Cada componente é o produto escalar com o respectivo unitário. Como o ângulo entre $\vet{K}$ e $\vet{k}$ é $\theta$ e o ângulo entre $\vet{K}$ e $\vet{i}$ é $90^\circ+\theta$,
\begin{align}
\vet{K}\cdot\vet{k} &= \cos\theta,\\
\vet{K}\cdot\vet{i} &= \cos(90^\circ+\theta) = -\seno\theta,
\end{align}
de onde \rel{K}
\begin{equation}
\vet{K}=-\seno\theta\,\vet{i}+\cos\theta\,\vet{k}.\label{m:K}
\end{equation}

\paragraph{Velocidade angular.} O corpo gira à taxa $\phid$ em torno de $Z$ e à taxa $\psid$ em torno de $z$, e as velocidades angulares se somam vetorialmente \rel{omega}. Substituindo a equação~\ref{m:K}:
\begin{align}
\vomega &= \phid\,\vet{K}+\psid\,\vet{k}
         = \phid\left(-\seno\theta\,\vet{i}+\cos\theta\,\vet{k}\right)+\psid\,\vet{k}\nonumber\\
        &= -\phid\seno\theta\,\vet{i}+\left(\psid+\phid\cos\theta\right)\vet{k}.
\end{align}
Define-se a componente axial \rel{compomega}
\begin{equation}
\omega_z=\psid+\phid\cos\theta.\label{m:wz}
\end{equation}

\paragraph{Velocidade angular do referencial.} O referencial $Oxyz$ acompanha somente a precessão: ele não gira com o \textit{spin}. Portanto \rel{Omega}
\begin{equation}
\vOmega=\phid\,\vet{K}=-\phid\seno\theta\,\vet{i}+\phid\cos\theta\,\vet{k}.\label{m:Omega}
\end{equation}

% ---------------------------------------------------------------------
\section{Quantidade de movimento angular}

Os eixos $Oxyz$ são principais de inércia em relação a $O$: por simetria, o tensor de inércia é diagonal e tem $I'$ nas duas direções transversais. Como $\omega_y=0$,
\begin{equation}
\vet{H}_O=\vet{I}_O\cdot\vomega=
\begin{pmatrix}I'&0&0\\0&I'&0\\0&0&I\end{pmatrix}
\begin{pmatrix}-\phid\seno\theta\\0\\\omega_z\end{pmatrix}
=\begin{pmatrix}-I'\phid\seno\theta\\0\\I\omega_z\end{pmatrix},
\end{equation}
isto é \rel{H},
\begin{equation}
\vet{H}_O=-I'\phid\seno\theta\,\vet{i}+I\omega_z\,\vet{k}.\label{m:H}
\end{equation}

\begin{figure}[H]
\centering
% Figura: composição de omega (a) e de H_O (b), em perspectiva, com a mesma vista da
% Figura 2(a). Todos os vetores pertencem ao plano vertical Zz. theta = 32 graus (esquemático).
\def\S{0.5299}
\def\C{0.8480}
\tikzset{
  vista3d/.style={scale=0.95,baseline=(cap.base),
    x={(0.966cm,0.129cm)},y={(-0.259cm,0.483cm)},z={(0cm,0.866cm)}},
  eixo/.style={-{Stealth[length=2.2mm]},tinta2,line width=0.5pt},
  vetor/.style={-{Stealth[length=2.6mm]},line width=1.3pt},
  resultante/.style={-{Stealth[length=3mm]},line width=1.7pt},
  aux/.style={tinta2,densely dashed,line width=0.4pt}}%
% #1: ângulo (em graus) em que o rótulo de theta é posicionado
\def\cenavet#1{%
  % eixo x: abaixo do plano de apoio, portanto desenhado antes dele
  \draw[eixo] (0,0,0) -- ({2.3*\C},0,{-2.3*\S}) coordinate (Xp);
  % plano de apoio horizontal
  \fill[apoio,opacity=0.8] (-1.8,-1.6,0) -- (2.6,-1.6,0) -- (2.6,2.2,0) -- (-1.8,2.2,0) -- cycle;
  \draw[tinta2,line width=0.3pt] (-1.8,-1.6,0) -- (2.6,-1.6,0) -- (2.6,2.2,0) -- (-1.8,2.2,0) -- cycle;
  \node[right,inner sep=2pt] at (Xp) {$x\ (\mathbf{i})$};
  % eixo y: atrás do plano Zz
  \draw[eixo] (0,0,0) -- (0,2.05,0) node[left,text=black,inner sep=1.5pt] {$y\ (\mathbf{j})$};
  % plano vertical Zz
  \fill[azul,opacity=0.10] (-1.6,0,0) -- (2.4,0,0) -- (2.4,0,4.3) -- (-1.6,0,4.3) -- cycle;
  \draw[azul!70,line width=0.3pt] (-1.6,0,0) -- (-1.6,0,4.3) -- (2.4,0,4.3) -- (2.4,0,0);
  \node[anchor=north west,font=\scriptsize,inner sep=2pt] at (-1.6,0,4.3) {plano $Zz$};
  % eixos no plano Zz
  \draw[eixo] (0,0,0) -- (0,0,4.85) node[above,text=black,inner sep=2pt] {$Z\ (\mathbf{K})$};
  \draw[eixo] (0,0,0) -- ({5.0*\S},0,{5.0*\C}) node[right,text=black,inner sep=2pt] {$z\ (\mathbf{k})$};
  % theta
  \draw[line width=0.5pt] plot[domain=0:32,samples=16] ({1.0*sin(\x)},0,{1.0*cos(\x)});
  \node[inner sep=1pt] at ({1.3*sin(#1)},0,{1.3*cos(#1)}) {$\theta$};
  \fill (0,0,0) circle (1.5pt) node[below left,inner sep=2pt] {$O$};}%
\begin{tikzpicture}[vista3d]
  \cenavet{9}
  \coordinate (A) at (0,0,2.15);
  \coordinate (B) at ({2.3*\S},0,{2.3*\C});
  \coordinate (W) at ({2.3*\S},0,{2.15+2.3*\C});
  \draw[aux] (A) -- (W);
  \draw[aux] (B) -- (W);
  \draw[vetor,azul] (0,0,0) -- (A);
  \draw[vetor,laranja] (0,0,0) -- (B);
  \draw[resultante] (0,0,0) -- (W);
  \node[left,inner sep=3pt] at (A) {$\dot{\phi}\,\mathbf{K}$};
  \node[below right,inner sep=2pt] at (B) {$\dot{\psi}\,\mathbf{k}$};
  \node[above left,inner sep=2pt] at (W) {$\boldsymbol{\omega}$};
  \node[anchor=north] (cap) at (canvas cs:x=0.45cm,y=-1.25cm) {\small (a) Velocidade angular};
\end{tikzpicture}\hfill
\begin{tikzpicture}[vista3d]
  \cenavet{21}
  \coordinate (Hz) at ({2.9*\S},0,{2.9*\C});
  \coordinate (Hx) at ({-1.2*\C},0,{1.2*\S});
  \coordinate (H) at ({2.9*\S-1.2*\C},0,{2.9*\C+1.2*\S});
  \draw[aux] (Hz) -- (H);
  \draw[aux] (Hx) -- (H);
  \draw[vetor,laranja] (0,0,0) -- (Hz);
  \draw[vetor,azul] (0,0,0) -- (Hx);
  \draw[resultante] (0,0,0) -- (H);
  \node[below right,inner sep=2pt] at (Hz) {$I\omega_z\,\mathbf{k}$};
  \node[left,inner sep=2pt] at (Hx) {$-I'\dot{\phi}\operatorname{sen}\theta\,\mathbf{i}$};
  \node[above right,inner sep=2pt] at (H) {$\mathbf{H}_O$};
  \node[anchor=north] (cap) at (canvas cs:x=0.45cm,y=-1.25cm) {\small (b) Quantidade de movimento angular};
\end{tikzpicture}

\caption{Composição vetorial de $\boldsymbol{\omega}$ e de $\mathbf{H}_O$ (a mesma figura do relatório).}
\label{m:fig:vet}
\end{figure}

A Figura~\ref{m:fig:vet} mostra que $\vet{H}_O$ tem direção diferente de $\vomega$: as componentes de $\vomega$ são multiplicadas por momentos de inércia diferentes.

% ---------------------------------------------------------------------
\section{Equação do movimento}

\paragraph{Derivada em referencial móvel.} Para um vetor qualquer, a derivada no referencial fixo é a derivada medida no referencial móvel mais o termo devido à rotação deste:
\begin{equation}
\dot{\vet{H}}_O=\left(\dot{\vet{H}}_O\right)_{Oxyz}+\vOmega\times\vet{H}_O.
\end{equation}
Na precessão permanente, $\theta$, $\phid$ e $\psid$ são constantes; logo, os coeficientes de $\vet{i}$ e $\vet{k}$ na equação~\ref{m:H} são constantes e a primeira parcela é nula. Assim \rel{eulerM},
\begin{equation}
\sum\vet{M}_O=\vOmega\times\vet{H}_O.\label{m:euler}
\end{equation}

\paragraph{Produto vetorial.} Com as equações~\ref{m:Omega} e~\ref{m:H}, escrevendo o produto como determinante:
\begin{equation}
\vOmega\times\vet{H}_O=
\begin{vmatrix}
\vet{i}&\vet{j}&\vet{k}\\
-\phid\seno\theta&0&\phid\cos\theta\\
-I'\phid\seno\theta&0&I\omega_z
\end{vmatrix}.
\end{equation}
Como a coluna central é nula a menos do unitário, as componentes $\vet{i}$ e $\vet{k}$ se anulam, e resta apenas a componente $\vet{j}$, que entra com o sinal negativo do cofator:
\begin{align}
\vOmega\times\vet{H}_O
&=-\vet{j}\left[\left(-\phid\seno\theta\right)\left(I\omega_z\right)-\left(\phid\cos\theta\right)\left(-I'\phid\seno\theta\right)\right]\nonumber\\
&=-\vet{j}\left[-I\omega_z\phid\seno\theta+I'\phid^{\,2}\seno\theta\cos\theta\right]\nonumber\\
&=\left(I\omega_z-I'\phid\cos\theta\right)\phid\seno\theta\,\vet{j}.\label{m:OxH}
\end{align}
Esse resultado é a equação~\ref*{rel-eq:OmegaH} do relatório.

\paragraph{Momento do peso.} O peso é aplicado em $G$, com $\vet{r}_G=c\,\vet{k}$. Usando a equação~\ref{m:K}, $\vet{W}=-W\vet{K}=W\seno\theta\,\vet{i}-W\cos\theta\,\vet{k}$. Lembrando que $\vet{k}\times\vet{k}=\vet{0}$ e $\vet{k}\times\vet{i}=\vet{j}$ \rel{M}:
\begin{equation}
\sum\vet{M}_O=c\,\vet{k}\times\left(W\seno\theta\,\vet{i}-W\cos\theta\,\vet{k}\right)=Wc\seno\theta\,\vet{j}.\label{m:M}
\end{equation}
A reação em $O$ não contribui, pois seu braço em relação a $O$ é nulo.

\paragraph{Da igualdade vetorial à equação escalar.} Substituindo as equações~\ref{m:OxH} e~\ref{m:M} na equação~\ref{m:euler}, ambos os lados têm a direção de $\vet{j}$. Para $\theta\neq0$, tem-se $\seno\theta\neq0$ e esse fator pode ser cancelado \rel{condicao}:
\begin{equation}
\phid\left(I\omega_z-I'\phid\cos\theta\right)=Wc.\label{m:cond}
\end{equation}

\paragraph{Forma de segundo grau.} Substitui-se $\omega_z$ pela equação~\ref{m:wz} e aplica-se a distributiva:
\begin{align}
\phid\left[I\left(\psid+\phid\cos\theta\right)-I'\phid\cos\theta\right]&=Wc\nonumber\\
I\psid\,\phid+I\cos\theta\,\phid^{\,2}-I'\cos\theta\,\phid^{\,2}&=Wc\nonumber\\
I\psid\,\phid-(I'-I)\cos\theta\,\phid^{\,2}&=Wc.\label{m:intermed}
\end{align}
Passando todos os termos para o lado direito, chega-se a \rel{quad}
\begin{equation}
(I'-I)\cos\theta\;\phid^{\,2}-I\psid\,\phid+Wc=0.\label{m:quad}
\end{equation}

\paragraph{Por que a massa não importa.} Como $I=mk^2$, $I'=mk'^2$ e $W=mg$, todos os coeficientes da equação~\ref{m:quad} contêm o fator $m$. Dividindo-a por $m$ \rel{quadk}:
\begin{equation}
(k'^2-k^2)\cos\theta\;\phid^{\,2}-k^2\psid\,\phid+gc=0.
\end{equation}

% ---------------------------------------------------------------------
\section{Item (a): as duas taxas de precessão}

Com $a=(I'-I)\cos\theta$ e $b=I\psid$, a equação~\ref{m:quad} é $a\phid^{\,2}-b\phid+Wc=0$. A fórmula de Bhaskara, com $\Delta=b^2-4aWc$, fornece \rel{raizes}:
\begin{equation}
\phid_1=\frac{b-\sqrt{\Delta}}{2a},\qquad\phid_2=\frac{b+\sqrt{\Delta}}{2a}.\label{m:raizes}
\end{equation}
Verificação dimensional: $a\phid^{\,2}$ tem unidade $\mathrm{kg\cdot m^2/s^2}=\mathrm{N\cdot m}$, a mesma de $Wc$ e de $b\phid$.

\subsection{Conta completa para $\theta=30^\circ$}

\begin{align}
a&=1,875\times10^{-4}\cdot\cos30^\circ=1,875\times10^{-4}\cdot0,866025=1,62380\times10^{-4}\un{kg\cdot m^2},\\
b^2&=(0,0117810)^2=1,38791\times10^{-4}\un{kg^2\cdot m^4/s^2},\\
4aWc&=4\cdot1,62380\times10^{-4}\cdot0,03924=2,54872\times10^{-5}\un{kg^2\cdot m^4/s^2},\\
\Delta&=1,38791\times10^{-4}-0,254872\times10^{-4}=1,13304\times10^{-4}\un{kg^2\cdot m^4/s^2},\\
\sqrt{\Delta}&=0,0106444\un{kg\cdot m^2/s},\\
2a&=3,24760\times10^{-4}\un{kg\cdot m^2}.
\end{align}
A raiz rápida resulta diretamente \rel{num30}:
\begin{equation}
\phid_2=\frac{0,0117810+0,0106444}{3,24760\times10^{-4}}=\frac{0,0224254}{3,24760\times10^{-4}}=69,05\un{rad/s}=659,4\un{rpm}.
\end{equation}

Na raiz lenta, o numerador é a diferença de dois números próximos, o que amplifica os erros de arredondamento. Multiplicando numerador e denominador da equação~\ref{m:raizes} por $b+\sqrt{\Delta}$ e usando $(b-\sqrt{\Delta})(b+\sqrt{\Delta})=b^2-\Delta=4aWc$, obtém-se uma forma equivalente, sem essa subtração:
\begin{equation}
\phid_1=\frac{4aWc}{2a\left(b+\sqrt{\Delta}\right)}=\frac{2Wc}{b+\sqrt{\Delta}}=\frac{2\cdot0,03924}{0,0224254}=3,4996\un{rad/s}=33,42\un{rpm}.\label{m:estavel}
\end{equation}

As duas raízes podem ser conferidas pelas relações de soma e produto de uma equação de segundo grau:
\begin{align}
\phid_1+\phid_2&=\frac{b}{a}=\frac{0,0117810}{1,62380\times10^{-4}}=72,55\un{rad/s}\quad(3,50+69,05=72,55),\\
\phid_1\,\phid_2&=\frac{Wc}{a}=\frac{0,03924}{1,62380\times10^{-4}}=241,7\un{rad^2/s^2}\quad(3,4996\cdot69,05=241,7).
\end{align}
Para $\theta<90^\circ$, $a$ é positivo; a soma e o produto são então positivos, e as duas raízes também: ambas as precessões têm o sentido do \textit{spin}.

\subsection{Toda a faixa de inclinação}

A Tabela~\ref{m:tab:a} repete a sequência acima para cada ângulo. Os coeficientes $b$ e $Wc$, e portanto $b^2=1,38791\times10^{-4}$, não dependem de $\theta$. A Tabela~\ref{m:tab:rpm} traz a conversão para rpm e as verificações.

\begin{table}[H]
\centering\small
\caption{Etapas do cálculo das raízes (unidades SI; $a$ em $10^{-4}$, $4aWc$ em $10^{-5}$, $\Delta$ em $10^{-4}$).}
\label{m:tab:a}
\begin{tabular}{cccccccc}
\toprule
$\theta$ & $\cos\theta$ & $a$ & $4aWc$ & $\Delta$ & $\sqrt{\Delta}$ & $\phid_1$ (rad/s) & $\phid_2$ (rad/s)\\
\midrule
$30^\circ$ & $0,86603$ & $1,6238$ & $2,5487$ & $1,1330$ & $0,010644$ & $3,4996$ & $69,052$\\
$35^\circ$ & $0,81915$ & $1,5359$ & $2,4108$ & $1,1468$ & $0,010709$ & $3,4895$ & $73,214$\\
$40^\circ$ & $0,76604$ & $1,4363$ & $2,2545$ & $1,1625$ & $0,010782$ & $3,4783$ & $78,543$\\
$45^\circ$ & $0,70711$ & $1,3258$ & $2,0810$ & $1,1798$ & $0,010862$ & $3,4660$ & $85,392$\\
$50^\circ$ & $0,64279$ & $1,2052$ & $1,8917$ & $1,1987$ & $0,010949$ & $3,4528$ & $94,296$\\
$55^\circ$ & $0,57358$ & $1,0755$ & $1,6880$ & $1,2191$ & $0,011041$ & $3,4387$ & $106,105$\\
$60^\circ$ & $0,50000$ & $0,9375$ & $1,4715$ & $1,2408$ & $0,011139$ & $3,4241$ & $122,240$\\
\bottomrule
\end{tabular}
\end{table}

\begin{table}[H]
\centering\small
\caption{Taxas em rpm e verificação da soma e do produto das raízes.}
\label{m:tab:rpm}
\begin{tabular}{ccccccc}
\toprule
$\theta$ & $\phid_1$ (rpm) & $\phid_2$ (rpm) & $\phid_1+\phid_2$ & $b/a$ & $\phid_1\phid_2$ & $Wc/a$\\
\midrule
$30^\circ$ & $33,42$ & $659,4$ & $72,55$ & $72,55$ & $241,7$ & $241,7$\\
$35^\circ$ & $33,32$ & $699,1$ & $76,70$ & $76,70$ & $255,5$ & $255,5$\\
$40^\circ$ & $33,22$ & $750,0$ & $82,02$ & $82,02$ & $273,2$ & $273,2$\\
$45^\circ$ & $33,10$ & $815,4$ & $88,86$ & $88,86$ & $296,0$ & $296,0$\\
$50^\circ$ & $32,97$ & $900,5$ & $97,75$ & $97,75$ & $325,6$ & $325,6$\\
$55^\circ$ & $32,84$ & $1013,2$ & $109,54$ & $109,54$ & $364,9$ & $364,9$\\
$60^\circ$ & $32,70$ & $1167,3$ & $125,66$ & $125,66$ & $418,6$ & $418,6$\\
\bottomrule
\end{tabular}
\end{table}

\subsection{\textit{Spin} mínimo}

Para que existam raízes reais, deve-se ter $\Delta\geq0$:
\begin{equation}
\left(I\psid\right)^2\geq4(I'-I)\cos\theta\;Wc
\quad\Longrightarrow\quad
\psid\geq\psid_{\mathrm{min}}=\frac{2}{I}\sqrt{(I'-I)\cos\theta\;Wc}.
\end{equation}
Essa é a equação~\ref*{rel-eq:psimin} do relatório. Numericamente, com $2/I=32000\un{kg^{-1}\cdot m^{-2}}$:
\begin{align}
\theta=30^\circ:&\quad\psid_{\mathrm{min}}=32000\sqrt{1,62380\times10^{-4}\cdot0,03924}=80,78\un{rad/s}=771\un{rpm},\\
\theta=60^\circ:&\quad\psid_{\mathrm{min}}=32000\sqrt{9,375\times10^{-5}\cdot0,03924}=61,38\un{rad/s}=586\un{rpm}.
\end{align}

% ---------------------------------------------------------------------
\section{Item (b): condição para $\psid\gg\phid$}

A equação~\ref{m:intermed} pode ser escrita, colocando $\phid$ em evidência, como
\begin{equation}
\phid\,\left(b-a\phid\right)=Wc.\label{m:bafi}
\end{equation}
A aproximação consiste em desprezar $a\phid$ em comparação com $b$, o que é razoável quando $\psid\gg\phid$. Resulta \rel{aprox}:
\begin{equation}
\phid_a=\frac{Wc}{b}=\frac{Wc}{I\psid}=\frac{mgc}{I\psid}=\frac{0,03924}{0,0117810}=3,3308\un{rad/s}=31,81\un{rpm}.
\end{equation}

\paragraph{Tamanho do termo desprezado.} Para $\theta=30^\circ$ e $\phid=\phid_1=3,4996\un{rad/s}$:
\begin{align}
I'\phid\cos\theta&=2,50\times10^{-4}\cdot3,4996\cdot0,866025=7,577\times10^{-4},\\
I\phid\cos\theta&=6,25\times10^{-5}\cdot3,4996\cdot0,866025=1,894\times10^{-4},\\
a\phid&=7,577\times10^{-4}-1,894\times10^{-4}=5,683\times10^{-4},
\end{align}
valor que corresponde a $5,683\times10^{-4}/0,0117810=4,82\%$ de $b$.

\paragraph{Por que a taxa não depende de $\theta$.} Admitir $\vet{H}_O\approx I\psid\,\vet{k}$ na equação~\ref{m:euler} equivale a escrever $\vOmega\times\vet{H}_O\approx\phid\,\vet{K}\times I\psid\,\vet{k}$, de módulo $\phid\,I\psid\seno\theta$. Igualando a $Wc\seno\theta$ \rel{OmegaH_approx}, o fator $\seno\theta$ se cancela.

\paragraph{Por que vale apenas para a raiz lenta.} Na raiz rápida, $a\phid_2$ não é pequeno diante de $b$: $\phid_2\approx b/a$, valor que é justamente o da soma das raízes (Tabela~\ref{m:tab:rpm}). Além disso, $\phid_2/\psid$ vai de $69,05/188,50=0,37$ a $122,24/188,50=0,65$, de modo que $\psid\gg\phid$ não é satisfeita.

% ---------------------------------------------------------------------
\section{Item (c): erro da aproximação}

\paragraph{Valor numérico.} Em $\theta=30^\circ$ \rel{erro30}:
\begin{equation}
e=\frac{\phid_a-\phid_1}{\phid_1}\times100\%=\frac{3,3308-3,4996}{3,4996}\times100\%=-4,82\%.
\end{equation}

\paragraph{Expressão exata.} Como $\phid_1$ é raiz, a equação~\ref{m:bafi} vale para ela. Dividindo-a por $b\,\phid_1$:
\begin{equation}
1-\frac{a\phid_1}{b}=\frac{Wc}{b\,\phid_1}=\frac{\phid_a}{\phid_1}.
\end{equation}
Como $e/100\%=\phid_a/\phid_1-1$, segue \rel{errofis}:
\begin{equation}
\frac{e}{100\%}=-\frac{a\phid_1}{b}.\label{m:e}
\end{equation}
Ou seja, o erro é exatamente a razão entre o termo desprezado e o termo mantido. Conferindo em $\theta=30^\circ$: $-5,683\times10^{-4}/0,0117810=-0,0482$.

\paragraph{Forma fechada.} Substitui-se $\phid_1$ da equação~\ref{m:raizes} na equação~\ref{m:e}; o fator $a$ se cancela e coloca-se $b$ em evidência na raiz (lembrando que $\sqrt{\Delta}=b\sqrt{1-4aWc/b^2}$):
\begin{equation}
\frac{a\phid_1}{b}=\frac{b-\sqrt{\Delta}}{2b}=\frac{1-\sqrt{1-4\eps}}{2},\qquad\eps=\frac{aWc}{b^2}.
\end{equation}
Logo \rel{errofechado}:
\begin{equation}
\frac{e}{100\%}=\frac{\sqrt{1-4\eps}-1}{2}.\label{m:efech}
\end{equation}
Para $\theta=30^\circ$: $\eps=1,62380\times10^{-4}\cdot0,03924/1,38791\times10^{-4}=0,04591$; $\sqrt{1-4\eps}=\sqrt{0,81636}=0,90353$; e $e=(0,90353-1)/2=-0,0482$.

\paragraph{Aproximação em série.} Pela série binomial $\sqrt{1-u}=1-u/2-u^2/8-\cdots$, com $u=4\eps$ \rel{erroserie}:
\begin{equation}
\sqrt{1-4\eps}=1-2\eps-2\eps^2-\cdots
\quad\Longrightarrow\quad
\frac{e}{100\%}=-\eps-\eps^2-\cdots.
\end{equation}
Em $\theta=30^\circ$, $-\eps$ fornece $-4,59\%$, e $-\eps-\eps^2$ fornece $-4,80\%$, próximo do valor exato de $-4,82\%$.

\paragraph{Dependência do \textit{spin}.} Como $b=I\psid$, tem-se $\eps\propto1/\psid^{\,2}$: ao dobrar o \textit{spin}, $\eps$ e, em primeira ordem, o erro caem a um quarto.

\begin{table}[H]
\centering\small
\caption{Etapas do cálculo do erro, com $\phid_a=3,3308\un{rad/s}$.}
\label{m:tab:erro}
\begin{tabular}{cccccc}
\toprule
$\theta$ & $\phid_1$ (rad/s) & $\eps$ & $\sqrt{1-4\eps}$ & $e$ exato (\%) & $-\eps$ (\%)\\
\midrule
$30^\circ$ & $3,4996$ & $0,04591$ & $0,90353$ & $-4,82$ & $-4,59$\\
$35^\circ$ & $3,4895$ & $0,04342$ & $0,90901$ & $-4,55$ & $-4,34$\\
$40^\circ$ & $3,4783$ & $0,04061$ & $0,91519$ & $-4,24$ & $-4,06$\\
$45^\circ$ & $3,4660$ & $0,03748$ & $0,92199$ & $-3,90$ & $-3,75$\\
$50^\circ$ & $3,4528$ & $0,03407$ & $0,92935$ & $-3,53$ & $-3,41$\\
$55^\circ$ & $3,4387$ & $0,03041$ & $0,93722$ & $-3,14$ & $-3,04$\\
$60^\circ$ & $3,4241$ & $0,02651$ & $0,94550$ & $-2,72$ & $-2,65$\\
\bottomrule
\end{tabular}
\end{table}

% ---------------------------------------------------------------------
\section{Reações no apoio e energia}

\paragraph{Reações.} O centro de gravidade descreve uma circunferência horizontal de raio $c\seno\theta$ com velocidade angular $\phid$. Sua aceleração é, portanto, centrípeta, de módulo $(c\seno\theta)\phid^{\,2}$. Aplicando a segunda lei de Newton \rel{reacoes}:
\begin{equation}
\text{vertical: } N-W=0\Rightarrow N=W;\qquad\text{horizontal: } R_h=mc\seno\theta\;\phid^{\,2}.
\end{equation}
A razão $R_h/W=c\seno\theta\,\phid^{\,2}/g$ não depende da massa. Em $\theta=30^\circ$, com $mc\seno\theta=0,100\cdot0,040\cdot0,5=0,0020\un{kg\cdot m}$:
\begin{align}
\text{lenta: }&R_h=0,0020\cdot(3,4996)^2=0,0245\un{N},&R_h/W&=0,025,\\
\text{rápida: }&R_h=0,0020\cdot(69,052)^2=9,536\un{N},&R_h/W&=9,72.
\end{align}

\paragraph{Energia.} A energia cinética tem uma parcela axial e uma transversal, e a energia potencial é medida a partir da altura de $O$ \rel{energia}:
\begin{equation}
E=\tfrac{1}{2}I'\left(\phid\seno\theta\right)^2+\tfrac{1}{2}I\omega_z^2+Wc\cos\theta.
\end{equation}
Em $\theta=30^\circ$, com $Wc\cos\theta=0,0340\un{J}$:
\begin{align}
\text{lenta: }&\omega_z=188,4956+3,4996\cdot0,866025=191,53\un{rad/s},\nonumber\\
&E=\tfrac{1}{2}(2,5\times10^{-4})(1,7498)^2+\tfrac{1}{2}(6,25\times10^{-5})(191,53)^2+0,0340=1,181\un{J},\\
\text{rápida: }&\omega_z=188,4956+69,052\cdot0,866025=248,30\un{rad/s},\nonumber\\
&E=\tfrac{1}{2}(2,5\times10^{-4})(34,526)^2+\tfrac{1}{2}(6,25\times10^{-5})(248,30)^2+0,0340=2,110\un{J}.
\end{align}
A Tabela~\ref{m:tab:fis} reúne os valores das três inclinações usadas no relatório.

\begin{table}[H]
\centering\small
\caption{Grandezas físicas das duas soluções.}
\label{m:tab:fis}
\begin{tabular}{ccccccc}
\toprule
$\theta$ & Solução & $\phid/\psid$ & $\omega_z$ (rad/s) & $R_h$ (N) & $R_h/W$ & $E$ (J)\\
\midrule
$30^\circ$ & lenta & $0,019$ & $191,53$ & $0,024$ & $0,025$ & $1,181$\\
 & rápida & $0,366$ & $248,30$ & $9,536$ & $9,721$ & $2,110$\\
$45^\circ$ & lenta & $0,018$ & $190,95$ & $0,034$ & $0,035$ & $1,168$\\
 & rápida & $0,453$ & $248,88$ & $20,624$ & $21,024$ & $2,419$\\
$60^\circ$ & lenta & $0,018$ & $190,21$ & $0,041$ & $0,041$ & $1,151$\\
 & rápida & $0,649$ & $249,62$ & $51,762$ & $52,765$ & $3,368$\\
\bottomrule
\end{tabular}
\end{table}

% ---------------------------------------------------------------------
\section{Como funciona a validação numérica}

A simulação resolve o movimento completo sem usar a equação~\ref{m:quad} nem os ângulos de Euler. Ela emprega eixos $1$, $2$ e $3$ solidários ao corpo, com o eixo $3$ sendo o de simetria. Esses eixos giram também com o \textit{spin}; não são os eixos $Oxyz$ da dedução.

\paragraph{Equações de Euler.} A equação vetorial do relatório (equação~\ref*{rel-eq:eulercorpo}), em componentes nos eixos do corpo, é
\begin{align}
I'\dot{\omega}_1&=M_1-(I-I')\,\omega_2\omega_3,\\
I'\dot{\omega}_2&=M_2+(I-I')\,\omega_1\omega_3,\\
I\dot{\omega}_3&=M_3,
\end{align}
em que $M_1$, $M_2$ e $M_3$ são as componentes do momento do peso nos eixos do corpo.

\paragraph{Orientação.} A orientação é representada por um quatérnio unitário $\vet{q}$, atualizado pela equação~\ref*{rel-eq:quat} do relatório. A partir dele obtém-se a matriz de rotação $\vet{R}$ e, com ela, o momento do peso em eixos do corpo:
\begin{equation}
\vet{M}=c\,\vet{e}_3\times\left(\vet{R}^{\mathsf{T}}\left(-W\vet{K}\right)\right),
\end{equation}
em que $\vet{e}_3$ é o vetor unitário do eixo $3$.

\paragraph{Condições iniciais.} O pião parte com inclinação $\theta_0$ e velocidade angular $\phid_0\vet{K}+\psid\vet{k}$, em que $\phid_0$ é a raiz que se deseja testar.

\paragraph{Integração.} O estado $\vet{s}=(\vet{q},\vomega)$ avança no tempo pelo método de Runge--Kutta de quarta ordem, com passo $h=2\times10^{-5}\un{s}$ e duração de $1\un{s}$, sendo $\vet{f}$ a função que dá $\dot{\vet{s}}$:
\begin{align}
\vet{k}_1&=\vet{f}(\vet{s}_n),\quad
\vet{k}_2=\vet{f}\left(\vet{s}_n+\tfrac{h}{2}\vet{k}_1\right),\quad
\vet{k}_3=\vet{f}\left(\vet{s}_n+\tfrac{h}{2}\vet{k}_2\right),\quad
\vet{k}_4=\vet{f}\left(\vet{s}_n+h\vet{k}_3\right),\\
\vet{s}_{n+1}&=\vet{s}_n+\tfrac{h}{6}\left(\vet{k}_1+2\vet{k}_2+2\vet{k}_3+\vet{k}_4\right).
\end{align}
Os vetores $\vet{k}_i$ acima são os estágios do método, sem relação com o unitário $\vet{k}$. A cada passo, $\vet{q}$ é renormalizado.

\paragraph{Leitura dos resultados.} O eixo de simetria no espaço é a terceira coluna de $\vet{R}$, com componentes $(\seno\theta\seno\phi,\,-\seno\theta\cos\phi,\,\cos\theta)$. Dele se extraem $\theta$ e $\phi$, e a taxa de precessão é o coeficiente angular da reta ajustada a $\phi(t)$. A Tabela~\ref{m:tab:val} compara os valores obtidos com as raízes analíticas.

\begin{table}[H]
\centering\small
\caption{Comparação entre as raízes analíticas e a integração numérica.}
\label{m:tab:val}
\begin{tabular}{cccc}
\toprule
$\theta$ & Solução & $\phid$ analítica (rad/s) & $\phid$ numérica (rad/s)\\
\midrule
$30^\circ$ & lenta & $3,4996$ & $3,4996$\\
           & rápida & $69,0524$ & $69,0524$\\
$45^\circ$ & lenta & $3,4660$ & $3,4660$\\
           & rápida & $85,3917$ & $85,3917$\\
$60^\circ$ & lenta & $3,4241$ & $3,4241$\\
           & rápida & $122,2396$ & $122,2396$\\
\bottomrule
\end{tabular}
\end{table}

\end{document}
